%%%%%%%%%%%%%%%%%%%%%%%%%%%%%%%%%%%%%%%%%
% Medium Length Professional CV
% LaTeX Template
% Version 2.0 (8/5/13)
%
% This template has been downloaded from:
% http://www.LaTeXTemplates.com
%
% Original author:
% Trey Hunner (http://www.treyhunner.com/)
%
% Important note:
% This template requires the resume.cls file to be in the same directory as the
% .tex file. The resume.cls file provides the resume style used for structuring the
% document.
%
%%%%%%%%%%%%%%%%%%%%%%%%%%%%%%%%%%%%%%%%%

%----------------------------------------------------------------------------------------
%	PACKAGES AND OTHER DOCUMENT CONFIGURATIONS
%----------------------------------------------------------------------------------------

\documentclass{resume} % Use the custom resume.cls style
\usepackage[left=0.75in,top=0.5in,right=0.75in,bottom=0.5in]{geometry} % Document margins
\usepackage[colorlinks = true,
            linkcolor = blue,
            urlcolor  = blue,
            citecolor = green,
            anchorcolor = blue]{hyperref}

%https://tex.stackexchange.com/questions/392703/change-one-footnote-symbol-from-number-to-asterisk


\usepackage{footnote}
\usepackage{footmisc}
\renewcommand{\thefootnote}{\alph{footnote}}
\newcommand{\astfootnote}[1]{%
\let\oldthefootnote=\thefootnote%
\setcounter{footnote}{0}%
\renewcommand{\thefootnote}{\fnsymbol{footnote}}%
\footnote{#1}%
\let\thefootnote=\oldthefootnote%
}




\newcommand{\tab}[1]{\hspace{.2667\textwidth}\rlap{#1}}
\newcommand{\itab}[1]{\hspace{0em}\rlap{#1}}

%\name{\begin{center}Curriculum Vitae\end{center}}
%\vspace{-2.00mm}
%\address{\begin{center}\large Siddhartha Gairola, Ph.D. Applicant in Computer Science and Engineering, U-M ID: 94142877\end{center}}

\name{\begin{center}Siddhartha Gairola\end{center}} % Your name
%\address{\begin{center}\small \textit{E-mail}: \href{mailto:siddhartha.gairola18@gmail.com}{sgairola@mpi-inf.mpg.de}, \textit{Website:} \href{https://sidgairo18.github.io/}{https://sidgairo18.github.io}\end{center}}% Your phone number and email
%\address{Microsoft Research Lab, 9, Vigyan 1st floor, Lavelle Road, Bengaluru, Karnataka 560001} % Your address
%\address{123 Pleasant Lane \\ City, State 12345} % Your secondary addess (optional)



%\address{\small\textit{Website:} \href{https://sidgairo18.github.io/}{https://sidgairo18.github.io}}
%\begin{small}\textit{E-mail: }siddhartha.gairola18@gmail.com\end{small}

\pagenumbering{arabic}

%\usepackage[english]{babel}
%\usepackage{fancyhdr}
%\usepackage{lastpage}

%\pagestyle{fancy}
%\fancyhf{}
%\rfoot{Page \thepage \hspace{1pt} of \pageref{LastPage}}

\begin{document}

%----------------------------------------------------------------------------------------
%	EDUCATION SECTION
%----------------------------------------------------------------------------------------

%\vspace{-2pt}
\begin{rSection}{Education}

{\bf MPI for Informatics \& Saarland University} \hfill {Saarbrücken, Germany}

Ph.D. Student, Computer Science \hfill {Oct, 2022 - present}
\begin{itemize} \itemsep -1.0mm
    \item Advisors: Prof. Bernt Schiele, Prof. Francesco Locatello
\end{itemize}

{\bf Institute of Science and Technology (ISTA)} \hfill {Klosterneuburg, Austria}

Visiting Researcher, Causal Learning and Artificial Intelligence Lab \hfill {Aug, 2024 - Feb, 2025}

{\bf International Institute of Information Technology, Hyderabad} \hfill {Hyderabad, India} 

MS by Research in CSE \hfill {2018 - 2020}
\begin{small}
\begin{itemize}\itemsep -1.0mm 
    \item Advisor: Prof. P. J. Narayanan
    \item Thesis: \textit{\href{https://web2py.iiit.ac.in/research_centres/publications/view_publication/mastersthesis/828}{Image Representations for Style Retrieval, Recognition and Background Replacement Tasks}}
    \item CGPA: \textbf{9.75/10}
\end{itemize}
\end{small}

BTech with Honours in CSE \hfill {2014 - 2018}
\begin{small}
\begin{itemize}\itemsep -1.0mm 
    \item CGPA: \textbf{8.91/10}
\end{itemize}
\end{small}


%{\bf St. Joseph's Academy} \hfill {Dehradun, India} 
%\\ Senior Secondary, ISC; Average: \textbf{93.6\%} \hfill {2012 - 2013}
%\\ Secondary, ICSE; Average: \textbf{95.28\%} \hfill {2010 - 2011}

\end{rSection}


%----------------------------------------------------------------------------------------
%	WORK EXPERIENCE SECTION
%----------------------------------------------------------------------------------------

%\vspace{-2pt}
\begin{rSection}{Work Experience} \itemsep -2pt

\begin{rSubsection}{\textbf{Microsoft Research, India} - \textit{Research Fellow}}{Aug 2020 - Aug 2022}{}{}
\begin{small}
\item Advisors: Dr. Mohit Jain, Dr. Nipun Kwatra
%\item Working on problems in healthcare, using computer vision and machine learning.
\item Developed a low-cost smartphone-based corneal topographer for keratoconus detection in collaboration with Sankara Eye Hospital. Published at ACM IMWUT/Ubicomp 2021.
\item Worked on training DNN models for abnormality detection in limited data settings. Published at EMBC 2021.
\end{small}
\end{rSubsection}

\vspace{-3pt}
\begin{rSubsection}{\textbf{Microsoft Research, India} - \textit{Research Intern}}{Jan 2020 - Jul 2020}{}{}
\begin{small}
\item Advisors: Dr. Mohit Jain, Dr. Nipun Kwatra
\item Worked on building smartphone based diagnostic solutions for eye diseases using mobile computing, computer vision and machine learning.
\end{small}
\end{rSubsection}

\vspace{-3pt}
\begin{rSubsection}{\textbf{Adobe Systems, India} - \textit{Product Intern}}{Jun 2019 - Jan 2020}{}{}
\begin{small}
\item Advisors: Balaji K., Mayur Hemani
\item Worked at the Media and Data Science Research lab.
\item Worked on Guided Few-shot Image Segmentation, improved on the prior state-of-the-art by 6\%. 
\item Published at IJCAI 2020 and filed a patent.
%\item Implemented a deep learning model in \textit{Python, PyTorch} to perform \textbf{Few-shot Guided Segmentation} on images, beat the SOTA by 6\%, published a paper (IJCAI 2020) and filed a patent for the same.
\end{small}
\end{rSubsection}

\vspace{-3pt}
\begin{rSubsection}{\textbf{IIIT Hyderabad} - \textit{Research Assistant}}{May 2016 - May 2019}{}{}
\begin{small}
\item Advisor: Prof. P. J. Narayanan
\item Research assistant at the Center for Visual Information and Technology (CVIT)
\item Worked on representation learning for image style search and retrieval. Published at WACV 2020.
\item Also worked on task of color-consistent background replacement. Published at MMM 2018.
\end{small}
\end{rSubsection}

\vspace{-3pt}
\begin{rSubsection}{\textbf{Google Summer of Code (GSoC)} - \textit{Student Developer}}{May 2018 - Aug 2018}{}{}
\begin{small}
\item Was accepted as a developer for the Google Summer of Code Program for the $2^{nd}$ time with Scilab.
\item Implemented a demo in C/C++ and Scilab as a working example for the MEX Library in Scilab.
\end{small}
\end{rSubsection}

\vspace{-3pt}
\begin{rSubsection}{\textbf{Google Summer of Code (GSoC)} - \textit{Student Developer}}{May 2017 - Aug 2017}{}{}
\begin{small}
\item Was accepted as a developer for the Google Summer of Code Program with the organization Scilab.
\item Implemented a C/C++ wrapper for Matlab MEX-API on current API Scilab.
\end{small}
\end{rSubsection}
\end{rSection}


% Patents
%----------------------------------------------------------------------------------------

\begin{rSection}{Patents}
%\vspace{-4pt}
\begin{rSubsection}{\large \underline{\textbf{}} }{}{}{}
 \item \textbf{S. Gairola}, M. Hemani, A. Chopra, J. Dahl, Balaji K. Improved Similarity Propagation for One-Shot and Few-Shot Image Segmentation \textbf{(US 16/906,954)}, \textit{2020}.
\end{rSubsection}
\end{rSection}



% Publications
%----------------------------------------------------------------------------------------
%\vspace{-7pt}
\begin{rSection}{Peer-Reviewed Publications} \itemsep -2pt
%\vspace{-2pt}
%\begin{rSubsection}{\large \underline {\textbf{In Submission}} }{}{}{}
%\item \textbf{S. Gairola}, M. Bohra, N. Shaheer, N. Jayashankar, P. %Joshi, A. Balasubramaniam, K. Murli, N. Kwatra, and M. Jain. SmartKC: %Smartphone-based Corneal Topographer for Keratoconus Detection. %\textit{In Proceedings of the ACM on Interactive, Mobile, Wearable and %Ubiquitous Technologies (IMWUT), Volume 5, Issue 4, 2021}.
%\end{rSubsection}

%\begin{rSubsection}{\large \underline{\textbf{Peer-Reviewed %Publications}} }{}{}{}

\item[(14.)] \textbf{S. Gairola*}, A. Pardyl*, S. Rao, A. Wróbel, B. Zieliński, B. Schiele, D. Rymarczyk. \href{https://}{DisParQ: Self-Supervised Part Concepts for Interpretable Vision Foundation Models.} \textit{Under Submission, 2026.}

\item[(13.)] \textbf{S. Gairola}, J. Xie, A. Kukleva, F. Locatello, B. Schiele. \href{https://}{Stranger Things: When Objects Appear Without Their Typical Neighbours.} \textit{In Conference on Neural Information Processing Systems (NeurIPS), 2026.}

\item[(12.)] \textbf{S. Gairola*}, A. Wróbel*, J. Tabor, B. Schiele, B. Zieliński, D. Rymarczyk. \href{https://arxiv.org/abs/2602.06613}{DAVE : Distribution-aware Attribution via ViT Gradient Decomposition.} \textit{In International Conference on Machine Learning (ICML), 2026.} \textbf{\textit{Spotlight, top 2.2\% of all $\sim$24K submissions.}}

\item[(11.)] \textbf{S. Gairola*}, R. Maser*, S. Rao, B. Schiele. \href{https://cvpr.thecvf.com/virtual/2026/poster/40400}{Align Once to Explain: Feature Alignment for Scalable B-cosification of Foundational Vision Transformers.} \textit{In Computer Vision and Pattern Recognition Conference (CVPR), 2026.}

\item[(10.)] \textbf{S. Gairola}, M. B{\"o}hle, F. Locatello, B. Schiele. \href{https://openreview.net/pdf?id=57NfyYxh5f}{How to Probe: Simple Yet Effective Techniques for Improving Post-hoc Explanations.} \textit{In International Conference on Learning Representations (ICLR), 2025.}

\item[(9.)] V. Ganatra, \textbf{S. Gairola}, P. Joshi, A. Balasubramaniam, K. Murali, A. Varadharajan, B. Mallikarjuna, N. Kwatra, M. Jain. \href{https://openaccess.thecvf.com/content/WACV2025/html/Ganatra_SmartKC_Improving_Performance_of_Smartphone-Based_Corneal_Topographers_WACV_2025_paper.html}{SmartKC++: Improving Performance of Smartphone-Based Corneal Topographers.} \textit{In IEEE Winter Conference on Applications of Computer Vision (WACV), 2025.}

\item[(8.)] \textbf{S. Gairola}, P. Joshi, A. Balasubramaniam, K. Murali, N. Kwatra, M. Jain. \href{https://arxiv.org/abs/2205.03702}{Keratoconus Classifier for Smartphone-based Corneal Topographer.} \textit{44th Annual International Conference of the IEEE Engineering in Medicine \& Biology Society (EMBC) (2022)}.

\item[(7.)] A. Aggarwal, \textbf{S. Gairola}, U. Uppadhyay, A. P. Vasishta, D. Rao, A. Goyal, K. Murali, N. Kwatra, and M. Jain. \href{https://arxiv.org/abs/2208.05552}{Towards Automating Retinoscopy for Refractive Error Diagnosis.} \textit{In Proceedings of the ACM on Interactive, Mobile, Wearable and Ubiquitous Technologies (IMWUT), Volume 6, Issue 3, 2022}.

\item[(6.)] \textbf{S. Gairola}, M. Bohra, N. Shaheer, N. Jayaprakash, P. Joshi, A. Balasubramaniam, K. Murali, N. Kwatra, and M. Jain. \href{https://sidgairo18.github.io/smartkc.pdf}{SmartKC: Smartphone-based Corneal Topographer for Keratoconus Detection.} \textit{In Proceedings of the ACM on Interactive, Mobile, Wearable and Ubiquitous Technologies (IMWUT), Volume 5, Issue 4, 2021}.

\item[(5.)] \textbf{S. Gairola}, F. Tom, N. Kwatra, and M. Jain. \href{https://arxiv.org/pdf/2011.00196.pdf}{RespireNet: A Deep Neural Network for Accurately Detecting Abnormal Lung Sounds in Limited Data Setting}. \textit{43rd Annual International Conference of the IEEE Engineering in Medicine and Biology Society (EMBC), 2021.}

\vspace{4pt}
\item[(4.)] \textbf{S. Gairola}, A. Chopra, M. Hemani, Balaji K. \href{https://www.ijcai.org/Proceedings/2020/0080.pdf}{SimPropNet:  Improved Similarity Propagation for Few-Shot Image Segmentation}. \textit{International Joint Conference on Artificial Intelligence (IJCAI), 2020}.

\vspace{4pt}
\item[(3.)] \textbf{S. Gairola}, R. Shah, P. J. Narayanan. \href{https://openaccess.thecvf.com/content_WACV_2020/papers/Gairola_Unsupervised_Image_Style_Embeddings_for_Retrieval_and_Recognition_Tasks_WACV_2020_paper.pdf}{Unsupervised Image Style Embeddings for Retrieval and Recognition Tasks}. \textit{IEEE Winter Conference on Applications of Computer Vision (WACV), 2020}. %\href{https://openaccess.thecvf.com/content_WACV_2020/html/Gairola_Unsupervised_Image_Style_Embeddings_for_Retrieval_and_Recognition_Tasks_WACV_2020_paper.html}{(link)}.

\vspace{4pt}
\item[(2.)] V. Kumar*, D. Khattar*, \textbf{S. Gairola}*, Y. K. Lal*, V. Varma. \href{https://arxiv.org/pdf/1710.01507.pdf}{Identifying Clickbait:  A Multi-Strategy Approach Using Neural Networks}. \textit{The 41st International ACM SIGIR Conference on Research \& Development in Information Retrieval, 2018}. %\href{https://dl.acm.org/doi/10.1145/3209978.3210144}{(link)}.

\vspace{4pt}
\item[(1.)] S. Rawat*, \textbf{S. Gairola*}, R. Shah, P. J. Narayanan. \href{https://sidgairo18.github.io/sky.pdf}{Find Me a Sky: A Data-Driven Method for Color-Consistent Sky Search and Replacement}. \textit{International Conference on Multimedia Modeling (MMM), 2018}. %\href{https://link.springer.com/chapter/10.1007/978-3-319-73603-7_18}{(link)}.
%\end{rSubsection}

%\vspace{-2pt}
\hfill \small{*equal contribution}
%\hfill \tiny{C=Conference, J=Journal}
% \astfootnote{equal contribution}
\end{rSection}

% Under Review
%----------------------------------------------------------------------------------------
%\vspace{-7pt}
% \begin{rSection}{Papers Under-Review} \itemsep -2pt
% %\vspace{-2pt}
% %\begin{rSubsection}{\large \underline {\textbf{In Submission}} }{}{}{}
% %\item \textbf{S. Gairola}, M. Bohra, N. Shaheer, N. Jayashankar, P. %Joshi, A. Balasubramaniam, K. Murli, N. Kwatra, and M. Jain. SmartKC: %Smartphone-based Corneal Topographer for Keratoconus Detection. %\textit{In Proceedings of the ACM on Interactive, Mobile, Wearable and %Ubiquitous Technologies (IMWUT), Volume 5, Issue 4, 2021}.
% %\end{rSubsection}

% %\begin{rSubsection}{\large \underline{\textbf{Peer-Reviewed %Publications}} }{}{}{}

% \item[(1.)] \textbf{S. Gairola}, M. B{\"o}hle, F. Locatello, B. Schiele. \href{https://}{How to Probe: Simple Yet Effective Techniques for Improving Post-hoc Explanations.} \textit{(under review)}.
% \end{rSection}

%	Academic Service
%----------------------------------------------------------------------------------------
%\vspace{-5pt}
\begin{rSection}{Academic Services} \itemsep -5pt
%\vspace{-6pt}
\item \textbf{Organizer:} Computer Vision for Developing Countries Workshop at \href{https://cv4dc.github.io/2024/}{ACCV 2024}, \href{https://cv4dc.github.io/2025/}{ICCV 2025}
\item \textbf{Reviewer:} ICML 2023-26 (\textbf{\textit{Gold reviewer in 2026}}); NeurIPS 2022-23,'25; ICLR 2022-24; CVPR 2024-26; ECCV 2024; ICCV 2025; IHCI 2021
\item \textbf{Volunteer:} ICVGIP 2018
\end{rSection}

%\clearpage
% Invited Talks and Presentations
%---------------------------------------------------------------------------------------
\vspace{-4pt}
\begin{rSection}{Talks \& Conference Presentations} \itemsep -4pt
\vspace{-4pt}
\begin{small}
\item Intriguing Applications and Overlooked Pitfalls of XAI in Visual Models, at the \href{https://gmum.net/}{GMUM  Workshop, Jagiellonian University}, October 2024.
\item A DNN for Accurately Detecting Abnormal Lung Sounds in Limited Data Setting, \textit{EMBC 2021} (\href{https://youtu.be/XoAF3fCAqq4}{video}).
\item Improved Similarity Propagation for Few-shot Image Segmentation, \textit{IJCAI 2020} (\href{https://www.ijcai.org/proceedings/2020/video/24830}{video}).
\item Unsupervised Image Style Embeddings for Retrieval and Recognition Tasks, \textit{WACV 2020} (\href{https://www.youtube.com/watch?v=EbUOg1gVcFw&t=909s}{video}).
\end{small}

\end{rSection}


%\clearpage
%	Projects
%----------------------------------------------------------------------------------------
%\vspace{-3pt}
\begin{rSection}{Selected Research Projects} %\itemsep -2.5pt
%\item \textbf{Hostel Portal:} Developed a hostel management portal using python, Web2py framework. \textit{Spring 2015}

%\item \textbf{Linux C-Shell:} A linux like terminal implemented using \textit{C} language using different concepts like fork, signals, exec, piping. \textit{Monsoon 2015}

%\item \textbf{P2P File Sharing:} Developed an application layer file transfer protocol with support for TCP. It included features such as file upload, file download and indexed searching. (Computer Networks course) \textit{Spring 2016}

\vspace{-8pt}

\item \textbf{Smartphone-based Corneal Topographer} \hfill Nov 2020 - Oct 2021 \\ 
Developed a low-cost smartphone-based corneal topographer. It provides a portable and scalable solution for the mass screening of keratoconus. Built the entire system: an AI powered smartphone app, the 3D printed attachment, and the processing software. [\textit{Android, Python, PyTorch}; \href{https://www.microsoft.com/en-us/research/project/smartkc-a-smartphone-based-corneal-topographer/}{project page}]


\item \textbf{Abnormality Detection in Respiratory Signals} \hfill Aug 2020 - Nov 2020\\ Developed a `ResNet34' based model along with a suite of novel augmentation and fine-tuning techniques to effectively utilize small-sized datasets. The method improved upon the state-of-the-art for 4-class classification on the ICBHI dataset by 2.2\%. [\textit{Python, PyTorch}; \href{https://github.com/microsoft/RespireNet}{code}] %\textit{2020}

\item \textbf{Similarity Propagation for Few-Shot Image Segmentation} \hfill Jun 2019 - Jan 2020\\
Developed a deep learning system that leverages background-foreground similarity to perform few-shot image segmentation. The method achieved state-of-the art results on the PASCAL-5i, COCO-20i and FSS datasets. [\textit{Python, PyTorch}] %\textit{Monsoon 2019}

\item \textbf{Unsupervised Image Style Representation Learning} \hfill Jan 2019 - Jun 2019\\
Developed a deep learning model to learn robust style representations without supervision. Achieved state-of-the-art results across six datasets and it was published as a paper. [\textit{Python, PyTorch}; \href{https://github.com/sidgairo18/unsupervised-style-learning}{code}] %\textit{Spring 2019}

\item \textbf{Sketching with Style} \hfill Aug 2018 - Dec 2018\\
Implemented the ICCV 2017 paper ``Sketching with Style'' (\href{https://openaccess.thecvf.com/content_iccv_2017/html/Collomosse_Sketching_With_Style_ICCV_2017_paper.html}{link}). Was able to reproduce the results provided in the paper and open-sourced the code.
%\footnote{J.  Collomosse,  T.  Bui,  M.  Wilber,  C.  Fang,  and  H.  Jin. Sketching with style:  Visual search with sketches and aesthetic context. In 2017 IEEE International Conference on Computer Vision (ICCV), 2017.} 
[\textit{PyTorch, Python}; \href{https://github.com/sidgairo18/Sketching-with-Style-ICCV-2017-PyTorch-}{code}]  %\textit{Monsoon 2018}

%\item \textbf{Movie Genre Classification} \hfill Jan 2018 - Apr 2018\\ 
%Developed a deep learning model to predict the genres of a movie based on its poster and plot. Created a new movie poster dataset by scraping IMDB website for movies produced in the US and UK. Identified four key components in a movie poster: key graphic, background, title text, billing block. Annotated the dataset with these four components using an OCR and saliency model. [\textit{Python, TensorFlow}] %\textit{Monsoon 2017}

% \item \textbf{Wiki-Search Engine} \hfill Aug 2017 - Dec 2017\\
% Implemented an efficient and scalable search engine on Wikipedia dump in \textit{Python}. Retrieve top relevant documents based on input query. Handled field as well as phrase queries. Documents were re-ranked based on `tf-idf' measure. For fast retrieval used threading and multi-level indexing.
%\textit{Monsoon 2017}

\item \textbf{Neural Clickbait Detection Engine} \hfill May 2017 - Aug 2017\\
Developed a system to robustly identify clickbait social media posts. The system comprises of a siamese-network that fuses visual and textual information to learn a neural embedding used to detect clickbaits. [\textit{Python, Keras}; \href{https://github.com/vaibhav4595/Clickbait_Detection}{code}] %\textit{Spring 2017}

\item \textbf{Color Consistent Background Replacement} \hfill Jan 2017 - Apr 2017\\
Developed a data driven method for color-consistent sky search and replacement. A diverse set of skies consistent in color and illumination are retrieved from a curated dataset, and used to generate composites. The composites are re-ranked based on realism and diversity. [\textit{MATLAB}; \href{https://cvit.iiit.ac.in/research/projects/cvit-projects/findmeasky}{project page}] %\textit{Spring 2017}

%\item \textbf{High Dynamic Range Images}: HDR imaging - generating images with a greater range of luminance levels than which can be achieved by taking only a single photograph with a fixed exposure. Using tone mapping, high-boost filtering and bilateral filtering for improvements.
%\textit{Monsoon 2016}

%\item \textbf{Tic-tac-toe Bot:} Designed an automated agent bot to play the a modified version of tic-tac-toe game. Using a minimax algorithm and alpha-beta pruning. \textit{Spring 2016}

%\item \textbf{2D Game:} A 2-dimensional shooting game developed using \textit{OpenGl C++}.\textit{Spring 2016}

%\item \textbf{3D Game:} A 3-dimensional game developed using \textit{OpenGl C++}(GLFW library).\textit{Spring 2016}

\end{rSection}

%\clearpage
% Teaching Experience
%---------------------------------------------------------------------------------------
\vspace{-5pt}
\begin{rSection}{Teaching Experience} \itemsep -3pt
\vspace{-5pt}
\item \textbf{Saarland University}: Elements of Data Science and Artificial Intelligence (Winter 2023, 2024)
\item \textbf{IIIT Hyderabad}: Computer Graphics (Spring 2019), Digitial Image Processing (Monsoon 2017, 2018), Computer Vision (Spring 2018), Artificial Intelligence (Spring 2017), Digital Logic and Processors (Monsoon 2016)
% \item \textbf{Computer Graphics}, IIIT Hyderabad \hfill Spring 2019
% \item \textbf{Digital Image Processing}, IIIT Hyderabad \hfill Monsoon 2018
% \item \textbf{Computer Vision}, IIIT Hyderabad \hfill Spring 2018
% \item \textbf{Digital Image Processing}, IIIT Hyderabad \hfill Monsoon 2017
% \item \textbf{Artificial Intelligence}, IIIT Hyderabad \hfill Spring 2017
% \item \textbf{Digital Logic and Processors}, IIIT Hyderabad \hfill Monsoon 2016

%\begin{rSubsection}{\textbf{CSE251: Computer Graphics}, IIIT Hyderabad}{Spring 2019}{}{}
%\item Head Teaching Assistant with Prof. P. J. Narayanan.
%\end{rSubsection}

%\vspace{-3pt}
%\begin{rSubsection}{\textbf{CSE478: Digital Image Processing}, IIIT Hyderabad}{Monsoon 2018}{}{}
%\item Teaching Assistant with Prof. Ravi Kiran.
%\end{rSubsection}

%\vspace{-3pt}
%\begin{rSubsection}{\textbf{CSE578: Computer Vision}, IIIT Hyderabad}{Spring 2018}{}{}
%\item Teaching Assistant with Prof. Avinash Sharma.
%\end{rSubsection}

%\vspace{-3pt}
%\begin{rSubsection}{\textbf{CSE478: Digital Image Processing}, IIIT Hyderabad}{Monsoon 2017}{}{}
%\item Teaching Assistant with Prof. Avinash Sharma.
%\end{rSubsection}

%\vspace{-3pt}
%\begin{rSubsection}{\textbf{CSE371: Artificial Intelligence}, IIIT Hyderabad}{Spring 2017}{}{}
%\item Teaching Assistant with Prof. Praveen Paruchuri.
%\end{rSubsection}

%\vspace{-3pt}
%\begin{rSubsection}{\textbf{IEC101: Digital Logic and Processors}, IIIT Hyderabad}{Monsoon 2016}{}{}
%\item Teaching Assistant with Prof. Madhav Krishna.
%\end{rSubsection}

The duties as a TA involved taking regular tutorials, setting up assignments and conducting evaluations.
\end{rSection}

%	Academic
%----------------------------------------------------------------------------------------
\vspace{-5pt}
\begin{rSection}{Awards \& Achievements}
\begin{rSubsection}{\large \underline {} }{}{}{} \itemsep -1.5pt
\item Awarded the \textbf{Dean's Research Award} for a paper published at an international conference while in undergraduate, 2018.
\item Awarded the \textbf{Dean's List Award} for 6 semesters straight: Monsoon'15, '16, '17 \& Spring'16, '17, '18
\item Qualified for the \textbf{ACM ICPC Asia Onsite Regionals} twice (2015, 2016).
%\item Was a member of the teams \textbf{Pendulum} and \textbf{GOAT} which qualified for the \textbf{ACM ICPC Asia Amritapuri Onsite Regionals} in 2015 and 2016 respectively.
%\item \textbf{All India Rank 2262} in \textbf{JEE Advanced 2014} (\textbf{98.5\%} percentile among 150,000 students).
%\item \textbf{All India Rank 3856} in \textbf{JEE Mains 2014} (\textbf{99.7\%} percentile among 1.2 million students).
\item \textbf{98.5} percentile in \textbf{JEE Advanced (2014)} examination among 150,000 students.
\item \textbf{99.7} percentile in \textbf{JEE Main (2014)} examination among 1.2 million students.
\end{rSubsection}
\end{rSection}



%----------------------------------------------------------------------------------------
%	TECHNICAL STRENGTHS SECTION
%----------------------------------------------------------------------------------------
%\vspace{-2.5pt}
\begin{rSection}{Technical Skills}
%\vspace{-1.0pt}
\begin{small}
\begin{tabular}{ @{} >{\bfseries}l @{\hspace{1ex}} l }
Working Knowledge &   Bash, C/C++, Python, MATLAB, Octave, OpenCV, Android-SDK \\
Software \& Tools & GNU/Linux, HTML, CSS, JavaScript, LaTeX, Excel \\
Deep Learning Frameworks & PyTorch, TensorFlow, Keras, Caffe \\
Past Experience & Java, OpenGL, WebGL, Django, Web2py 
\end{tabular}
\end{small}

\end{rSection}

%----------------------------------------------------------------------------------------
%	Contact Details
%----------------------------------------------------------------------------------------
%\vspace{-2.5pt}
\begin{rSection}{Contact}
%\vspace{-1.0pt}
\begin{small}
\begin{tabular}{ @{} >{\bfseries}l @{\hspace{1ex}} l }
E-mail &   \href{mailto:sgairola@mpi-inf.mpg.de}{sgairola@mpi-inf.mpg.de} \\
Website & \href{https://sidgairo18.github.io/}{https://sidgairo18.github.io/} \\
Address & Office 608, Max-Planck-Institut für Informatik, Campus E1 4, 66123 Saarbrücken, Germany \\
Phone & +49 681 9325 2149 \\
Fax & +49 681 9325 2099 \\
\end{tabular}
\end{small}

\end{rSection}


%	Programming and Open-source
%----------------------------------------------------------------------------------------
%\vspace{-3pt}
%\begin{rSection}{Programming and Open-Source} \itemsep -5pt
%\item \textbf{Programming:} Take part actively in competitive programming contests.\newline
%\textbf{Codechef} (handle: sidgairo18)\newline
%\textbf{Codeforces} (handle: sidgairo18)

%\item \textbf{Open-Source:} Contribute to open source actively. Have made contributions to open source organizations like \textbf{Scilab, LibreOffice and CCExtractor}.
%(\textbf{Github profile} : sidgairo18)
%\end{rSection}

%----------------------------------------------------------------------------------------
%\vspace{-3pt}
%\begin{rSection}{Advanced Courses}
%\itab{\textbf{Core Courses}} \tab{}  %\tab{\textbf{Other Courses}}
%\\ \itab{Computer Programming } \tab{}  \tab{Linear Algebra}
%\\ \itab{Data Structures} \tab{}  \tab{Discrete Mathematics} 
%\\ \itab{Algorithms} \tab{}  \tab{Complexity and Advanced Algorithms} 
%\\ \itab{Operating Systems} \tab{} \tab{Computer System \& Organization}
%\\ \itab{Database Systems} \tab{} \tab{Digital Image Processing}
%\\ \itab{Artificial Intelligence} \tab{} \tab{Computer Vision}
%\\ \itab{Machine Learning} \tab{} \tab{Information Security}
%\\ \itab{Advanced Computer Networks} \tab{} \tab{Systems Network Security}
%\\ \itab{Information Retrieval \& Extraction} \tab{} \tab{Natural Language Applications}
% \\ \itab{Process Control (ongoing)} \tab{} \tab{Electrodynamics}

%\end{rSection}

%----------------------------------------------------------------------------------------
%\begin{rSection}{Extra-Cirrucular} \itemsep -5pt
%\item Play the piano and guitar.
%\item Member of the school football team from 2007-2011.
%\item Member of the IIIT Hyderabad college football team from 2014-2018.
%\end{rSection}


%	References
%----------------------------------------------------------------------------------------
%\vspace{-4pt}
%\begin{rSection}{References} \itemsep -5pt
%\vspace{-4pt}
%\begin{small}
%\item \textbf{Prof. P. J. Narayanan}, Director, IIIT-Hyderabad, India \href{mailto:pjn@iiit.ac.in}{(email)}
%\item \textbf{Dr. Mohit Jain}, Senior Researcher, Microsoft Research, India \href{mailto:mohja@microsoft.com}{(email)}
%\item \textbf{Dr. Nipun Kwatra}, Principal Researcher, Microsoft Research, India \href{mailto:nkwatra@microsoft.com}{(email)}
%\item \textbf{Balaji Krishamurthy}, Principal Scientist, Adobe Systems, India \href{mailto:kbalaji@adobe.com}{(email)}
%\item \textbf{Mayur Hemani}, Senior Research Scientist, Adobe Systems, India \href{mailto:mayur@adobe.com}{(email)}
%\end{small}
%\end{rSection}



\end{document}
